\section{Gauge comparison and the geometry of large defects}
\label{sec:seed-geometry}
We now prove the geometric estimates used in the seed energy bound.
The essential analytic point is a first-power estimate for the actual
coarse predictor. It must hold on both branches of its definition and
for every output with positive weight. Small-field uniqueness alone
would not give the estimate needed here.

Write $r(g)=D^{-1}\|g-1\|$ for the normalized matrix chord. Thus
$r\le2$, $r(g^{-1})=r(g)$, $r$ is conjugation invariant, and
$r(gh)\le r(g)+r(h)$. For a finite lattice box $X$, let
\[
 E_X(U)=\sum_{\substack{(x,a)\\x,x+e_a\in X}}r(U_{x,a}),\qquad
 F_X(U)=\sum_{\substack{(x,p)\\\text{all corners in }X}}r(U_p).
\]
The sums include positive bonds and one orientation for each plaquette.
All boxes in the first part of the argument lie in the covering lattice
$\mathbb Z^4$; periodic multiplicities will be restored at the end.

\subsection{An axial gauge with a first-power bound}
For the box of radius $R$, with side $q=2R+1$, there is a coordinate
path gauge $H$ such that
\begin{equation}\label{eq:geometry-axial}
 E_X(H\cdot U)\le16q F_X(U).
\end{equation}
Here the gauge at $x$ is parallel transport from a suitable base point
to $x$, moving along the coordinates in a prescribed order. We prove
more generally the bound $4kqF_X$ on a $k$-dimensional coordinate slab.
It is zero on a one-point slab. Suppose the assertion holds in $k$
dimensions and add a coordinate $j$. Choose a slice perpendicular to
$j$ with least plaquette sum; its sum is at most $F_X/q$, since the
plaquettes in different slices are distinct. Choose the base point
there by induction, and extend the paths by moving last along $j$.
All $j$-bonds in the resulting gauge are the identity.

For a transverse bond, two neighbouring slice values differ in chord
by at most the chord of the plaquette between them. Indeed, with the
two vertical bonds equal to $1$, that plaquette is $ab^{-1}$, and
$|r(a)-r(b)|\le r(ab^{-1})$. Summing over the transverse directions
costs at most four times the plaquette sum between the slices.
Telescoping in $j$ shows that the bond sums of any two slices differ
by at most $4F_X$. The total bond sum is consequently at most
\[
 q\left(4kq\frac{F_X}{q}+4F_X\right)=4(k+1)qF_X.
\]
This proves \eqref{eq:geometry-axial}, with the required coordinate
path form of the gauge. In particular the fallback rule which permits
axial gauges with bound $K_{\rm ax}qF_X$ has an eligible choice whenever
$K_{\rm ax}\ge16$.

We also need a path counting estimate. If $N=|X|$, then for any fixed
coordinate order,
\begin{equation}\label{eq:geometry-path-count}
 \sum_{x,y\in X}r(U_{x\to y})\le4Nq E_X(U).
\end{equation}
For each coordinate segment use the sum of the bond chords on its
coordinate line. During that segment the other coordinates are taken
from either $x$ or $y$, according to whether they have already been
changed. Exchanging these coordinate slots between the two endpoints
is a bijection of $X\times X$. Thus the mixed point is uniformly
distributed when the endpoints are uniformly distributed. Averaging
the sum over its coordinate line gives $qE_{X,a}/N$, where $E_{X,a}$
is the bond sum in direction $a$. Bounding each of the at most four
segments by $qE_X/N$ and multiplying by $N^2$ proves
\eqref{eq:geometry-path-count}. Reversed segments have the same bound
because the chord is invariant under inversion.

\subsection{Synchronizing minimizing frames}
Let $I$ and $A$ be nonempty finite sets, and suppose
$W_{ij\alpha}\in G$, with $W_{ii\alpha}=1$ and
$W_{ji\bar\alpha}=W_{ij\alpha}^{-1}$ for a permutation
$\alpha\mapsto\bar\alpha$ of $A$. A minimizing frame $h:I\to G$
minimizes
\[
 \mathcal E(h)=\mathop{\rm av}_{i,j,\alpha}
       \tfrac12 r(h_iW_{ij\alpha}h_j^{-1})^2.
\]
Such a minimum exists by compactness. Put
\[
 b_i=\mathop{\rm av}_{j,\alpha}r(W_{ij\alpha}),\qquad
 b=\mathop{\rm av}_i b_i.
\]
There is a constant $C_G$, independent of both finite sets and of $W$,
such that some $c\in G$ satisfies
\begin{equation}\label{eq:geometry-cluster}
 \mathop{\rm av}_i r(h_ic^{-1})\le C_G b.
\end{equation}

Here are the details which improve the square-root estimate to this
first-power bound. Comparison with the constant frame gives
$\mathcal E(h)\le b$, since $r^2/2\le r$. The triangle inequality and
$(u+v)^2\le2u^2+2v^2$ therefore give
\begin{equation}\label{eq:geometry-pair-square}
 \mathop{\rm av}_{i,j}\|h_i-h_j\|^2\le8D^2b.
\end{equation}
Write $m=\mathop{\rm av}_j h_j$ for the matrix centroid and
$m_i=\mathop{\rm av}_{j,\alpha}h_jW_{ij\alpha}^{-1}$ for its local
replacement. Then
\begin{equation}\label{eq:geometry-feedback}
 \|m_i-m\|\le Db_i,
 \qquad h_i\text{ is a nearest group element to }m_i.
\end{equation}
The norm bound follows by averaging matrix differences. For the second
assertion, change only $h_i$ in the minimizing energy. The two
 off-diagonal parts have equal costs by the reversal identity. The
 diagonal term in the actual energy remains zero. Freezing its second
 endpoint would add the nonnegative term
$r(kh_i^{-1})^2/2$, which is zero at $k=h_i$; hence $h_i$ also minimizes
the full frozen-neighbour cost. Expanding the squared Frobenius norms
shows that this is exactly nearest-point minimization for $m_i$.

The local tubular projection of the group has derivative the
orthogonal tangent projection at $1$. On a sufficiently small ambient
ball its derivative has norm at most $2$. Shrinking the ball also
ensures that every nearest point is this local projection. Thus choose
$\delta>0$ so that two matrices in the $\delta$-ball have nearest
points at distance at most twice their distance. This also follows
directly from the first-order nearest-point estimate used for the
group data. Set
\[
 \sigma=\frac{(\delta/2)^2}{8D^2},\qquad a=\frac{\delta}{2D}.
\]
If $b<\sigma$, choose $c$ nearest to $m$ and rotate the frame by
$c^{-1}$. In the rotated coordinates $1$ is nearest to its centroid.
Choose $k$ whose mean squared distance to the frame is at most the
double average in \eqref{eq:geometry-pair-square}. Jensen's inequality
and the nearest-point property give
\[
 \|m-1\|^2\le\|m-h_k\|^2
 \le\mathop{\rm av}_j\|h_j-h_k\|^2< (\delta/2)^2.
\]
For $b_i<a$, \eqref{eq:geometry-feedback} puts $m_i$ in the
$\delta$-ball and gives $Dr(h_i)\le2Db_i$. For $b_i\ge a$ use
$r(h_i)\le2\le2b_i/a$. Averaging proves the desired inequality in
this case. If $b\ge\sigma$, simply use the bound $2\le2b/\sigma$.
Thus \eqref{eq:geometry-cluster} holds with
$C_G=\max\{2,2/a,2/\sigma\}$. Left and right rotations give the same
chord, by conjugation invariance.

Apply this to the coordinate path transports on a patch $X$ of side
$q$. The path orders supply $A$, with reversed orders supplying the
involution above. If $f$ is the actual minimizing frame and $H$ is any
comparison gauge, then $h_i=f_iH_i^{-1}$ minimizes the path cost for
$H\cdot U$. This follows by substituting into every path product;
normalizing a competitor at the patch centre is harmless because a
common left rotation preserves the cost. Equations
\eqref{eq:geometry-path-count} and \eqref{eq:geometry-cluster} yield
\[
 \sum_{i\in X}r(c^{-1}f_iH_i^{-1})\le4C_Gq E_X(H\cdot U).
\]
For an active fallback frame, eligibility instead gives
$E_X(f\cdot U)\le K_{\rm ax}qF_X(U)$. The pure-gauge discrepancy
$d_i=f_iH_i^{-1}$ satisfies
\[
 E_X(d\cdot1)\le E_X(f\cdot U)+E_X(H\cdot U).
\]
This is the bondwise triangle inequality, with conjugation removing
$H\cdot U$. Its path transport from $i$ to $j$ is $d_id_j^{-1}$.
Average \eqref{eq:geometry-path-count} over $j$ and select one $j$
whose sum is no larger than the average. Taking $c=d_j$ proves the
same conclusion with the added axial term. Both branches therefore
satisfy, for $A_G=\max\{4C_G,4\}$,
\begin{equation}\label{eq:geometry-frame}
 \sum_{i\in X}r(c^{-1}f_iH_i^{-1})
 \le A_G q\{E_X(H\cdot U)+K_{\rm ax}qF_X(U)\}.
\end{equation}
Only positivity of the selected frame weight is required; the
fallback assertion would be false for an ineligible tag.

\subsection{The actual coarse predictor}
Fix the smoothing order $J\ge4$ and the kernel constants. The odd block
side $L$ may vary. The patch has radius $(J+3)L$, and the macro box
has radius $(J+7)L$. Fix a coarse site $x$, lift the field around $Lx$
to this macro box, and write $P_x$ for its plaquette chord sum.
All translated patches at coarse offsets in the radius-two box lie
inside the macro box. Their connecting straight paths of length $L$
lie there as well.

By \eqref{eq:geometry-axial} choose a macro gauge $H$ with
$E_{\rm macro}(H\cdot U)\le D_0 LP_x$, where $D_0$ depends only on
$J$. Applying \eqref{eq:geometry-frame} on each translated patch gives
constants $c_y\in G$ with
\begin{equation}\label{eq:geometry-macro-frame}
 \sum_i r(c_y^{-1}f_{x+y}(i)H_{i+Ly}^{-1})\le C_0L^2P_x.
\end{equation}
Indeed each patch side is at most a constant times $L$, and its bond
and plaquette sums are bounded by the macro sums. The constants
$C_0,D_0$ are fixed before choosing $L$.

The patch weights $w_i$ are the distribution of the sum of $J$
independent uniform points in the small box of side $L$. They obey
\begin{equation}\label{eq:geometry-light}
 w_i\le L^{-4}.
\end{equation}
Condition on all but the first point. For a prescribed sum there is
at most one possible first point, among $L^4$ equiprobable choices.
This proves \eqref{eq:geometry-light}, including coincidences among
different tuples of points.

For a coarse bond from $x+y$ to $x+y+e_a$, let $Y_i$ be its connector
in the two selected frames. Gauging the middle straight path by $H$
expresses $c_y^{-1}Y_ic_{y+e_a}$ as the product of the two frame
errors and that path transport. At each of its $L$ steps, translation
of $i$ is injective; hence the sum of the path chords is at most
$L E_{\rm macro}(H\cdot U)$. Equations
\eqref{eq:geometry-macro-frame} and \eqref{eq:geometry-light} give
\begin{equation}\label{eq:geometry-connector}
 \sum_i w_i r(c_y^{-1}Y_ic_{y+e_a})
       \le(2C_0+D_0)L^{-2}P_x.
\end{equation}

The predictor $Q$ either minimizes the weighted squared chord to the
connectors, or selects a connector when that minimum is at least the
fixed positive cutoff $\zeta$. For every group element $b$ its output
satisfies
\[
 r(Qb^{-1})\le2\max\{1,\zeta^{-1}\}\sum_iw_i r(Y_ib^{-1}).
\]
On the minimizing branch, $Q$ is nearest to the weighted matrix
centroid; comparison with $b$ and two triangle inequalities prove the
factor $2$. On the other branch,
$\zeta\le\sum_iw_i r(Y_ib^{-1})^2/2\le\sum_iw_i r(Y_ib^{-1})$,
while $r(Qb^{-1})\le2$. This proves the claim for every output tag.
Using $b=c_yc_{y+e_a}^{-1}$ in \eqref{eq:geometry-connector} bounds
each gauged coarse bond by $C_1L^{-2}P_x$. The four corner gauges
cancel around a plaquette, so
\begin{equation}\label{eq:geometry-predictor-curvature}
 r(Q_p)\le C L^{-2}P_x,\qquad
 r(V_p)\le\sum_{e\in\partial p}r(V_eQ_e^{-1})+C L^{-2}P_x,
 \quad C=4C_1.
\end{equation}
The second inequality follows by replacing the four factors of the
plaquette one at a time, conjugating the intervening products.
No smallness assumption on $U$ was used. The hypothesis is that the
frames at the radius-two coarse offsets have positive weight.

\subsection{Periodic multiplicity and the large-defect witness}
On a periodic lattice, let $m(x,p)$ count the lifts of the fine
plaquette $p$ wholly contained in the macro box centred at $Lx$.
Then, with $r_p=r(U_p)$,
\begin{equation}\label{eq:geometry-port}
 P_x=\sum_p m(x,p)r_p.
\end{equation}
This equality counts lifts, not merely distinct periodic plaquettes.
If the fine period is greater than $2(J+7)L$, different anchors in
one macro box cannot wrap to the same point, and $m(x,p)\le1$.
Furthermore, if $m(x,p)>0$, the coarse parent of $p$ lies at periodic
sup distance at most $R=J+8$ from $x$: the fine anchor is within
$(J+7)L$ of $Lx$, its distance to its block centre is at most $L/2$,
and distances between embedded coarse sites are $L$ times their
coarse distances. Consequently, under the same nonwrapping condition,
\begin{equation}\label{eq:geometry-port-column}
 \sum_xm(x,p)\le(2R+1)^4.
\end{equation}
The periodic ball bound holds even when the ball wraps, since reducing
an integer box modulo the period is surjective onto that ball.

A seed collection consists of fine plaquettes $\mathcal F$, coarse
predictor bonds $\mathcal I$ and coarse flags $\mathcal G$. Its size
is $n=|\mathcal F|+|\mathcal I|+|\mathcal G|$. Associate to it the
coarse parents of the fine plaquette corners, the endpoints of the
predictor bonds and the flag sites. Their union has size at most
$4n$. A radius-$b$ neighbourhood of this union therefore contains at
most
\[
 v(b)n,\qquad v(b)=4(2b+1)^4 N_p,
\]
coarse plaquette addresses, where $N_p$ is the number of plaquette
orientations. For $B\ge1$, $v(2B)\le v(2)B^4$. This proves the
outer-region bound used in Section~\ref{sec:seed-energy}.

Let $t$ be the fine threshold, $T$ the very large threshold, and $W$
the predictor cutoff parameter. Retain the seeds in $\mathcal F$
with $r_p>T$ and those in $\mathcal I$ with
$r(V_eQ_e^{-1})>T$, and discard the flags. Call this collection the
very large seeds, of size $d_A$. The local sector conditions say
that, in the radius-$5B$ seed neighbourhood, unlisted fine chords
are at most $t/2$, unlisted predictor chords at most $2Wt$, and all
selected frame weights are positive.

Suppose
\begin{equation}\label{eq:geometry-witness-conditions}
 R\le\lfloor4B/5\rfloor,\quad t/2\le T,\quad2Wt\le T,\quad
 \left(4+CL^{-2}(2(J+7)L+1)^4N_p\right)T\le\delta.
\end{equation}
Then every coarse plaquette of chord greater than $\delta$ whose
anchor lies in the radius-$2B$ seed neighbourhood is within distance
$R$ of a very large seed. To prove this, suppose no such seed exists.
Every fine plaquette counted by $P_x$ has parent within $R$ of $x$.
If listed, its chord is at most $T$, or it would be a witness; if
unlisted, its chord is at most $t/2\le T$. The same reasoning applies
to the four predictor discrepancies, using $2Wt\le T$. All these
addresses and the required radius-two frames lie in the radius-$5B$
neighbourhood, since $x$ starts within $2B$ and the added radii are
at most $3B$. Thus \eqref{eq:geometry-predictor-curvature} applies and
gives $r(V_p)\le\delta$, a contradiction.

It follows that these large coarse plaquettes belong to the hit
region, defined using the radius-$\lfloor4B/5\rfloor$ neighbourhood
of the very large seeds, and their number is at most $v(R)d_A$.
The whole hit region has at most $v(2)B^4d_A$ plaquettes. A coarse
flag is, by definition, a corner of a coarse plaquette with chord
above the coarse threshold $s$; each such plaquette has at most four
corners. Hence
\[
 |\mathcal G|\le4\#\{p\text{ in the outer region}:r(V_p)>s\}.
\]
Splitting these plaquettes at $\delta$ gives precisely the flag count
and the $K_0d_A$ large-plaquette count used in the energy allocation.

All constants in these geometric counts can be chosen before $L$
and $W$: take, for example,
$K_0=v(R)+v(2)+R+(2R+1)^4$. For each fixed $L,W$, the conditions
\eqref{eq:geometry-witness-conditions} hold eventually along the
coupling parameter. Indeed $B\to\infty$, $T=h^{-49/100}\to0$, and
$t/T=P(h)h^{-1/100}\to0$, by the threshold estimates already proved.
The allowed coarse periods are at least $100(1+B)$, which also
ensures the nonwrapping condition for the macro box. Thus the
geometric constants are independent of both the volume and the
subsequent choice of block side. The guide bounds and the hit and
unhit energy comparison estimates are separate analytic inputs;
this argument establishes the geometric part of their allocation.
